\chapter{Introduction}
%	\begin{tikzpicture}[remember picture,overlay]
	%\node[anchor=east,inner sep=0pt] at (current page text area.east|-0,3cm) %{\includegraphics[height=5.5cm]{writing_tool_cmp.png}};
%	\end{tikzpicture}



Scientists are writers by trade. We make our living by writing papers and other pieces of literature and yet interestingly most of us are never formally taught how to write. This is made worse by the fact that writing is a relatively nuanced field that isn't easy to figure out by feel. This document was motivated by the lack of standard resources we faced while writing our Master's thesis. As your Master's thesis is a significant portion of your college grade, it is likely to be the first piece of literature that you would need to and would want to be as meticulous as you can.

We, like many others, faced this exact same issue when we wrote our Master's thesis and we decided that we should do something about it. We had collected a lot of snippets and advice from our conversations with our guides and professors. This document seeks to collect all that into a single place to create a reference for when you write reports.

It is to be noted that this guide is \textbf{not an exhaustive reference.} It does not contain all the best practices and things to do to write the best report. It is written by students who, like most of you, were never formally taught to write. You can expect this document to improve your report (maybe even significantly), however, we do not claim to how to write a perfect report. There will definitely be things that your guide will help you with that will be missing from this guide and when that happens we strongly urge you to mail any of the authors of this document. We would love to add whatever knowledge you can contribute to this document.

\LaTeX\ might seem like a unnecessarily difficult choice for writing documents. While this might be true for something quite simple, we as scientists, rarely have the pleasure of writing such simple documents. Our work, even as undergrad students, is simply too complicated to be typeset\footnote{\url{https://en.wikipedia.org/wiki/Typesetting}} without proper tools. So while simpler documents might be easier in a more mainstream word processor such as LibreOffice (or god forbid MS Office), the more complex a document gets the more \LaTeX\ gets to shine. You can refer to Figure on the top of this page for a handy guide on when to use which tool for your writing exploits.

\begin{center}
	\includegraphics[width=0.75\linewidth]{images/writing_tool_cmp}
\end{center}
